\documentclass[12pt]{article}

\usepackage{url,verbatim, amsmath, amssymb, amsfonts, dsfont, color}
\textwidth 15true cm
\textheight 9true in
\oddsidemargin 0.30true in
\evensidemargin 0.30true in
\topmargin -0.3true in
\headsep 0.4true in

\def\sgn{{\rm sign}}
\def\drightarrow{{\stackrel{d}{\rightarrow}}}
\def\dotsim{\stackrel{.}{\sim}}
\def\goesd{\stackrel{d}{\rightarrow}}
\def\goesp{\stackrel{p}{\rightarrow}}
\def\goeslaw{\stackrel{{\cal L}}{\longrightarrow}}
%
\def\vp{{\varphi}}
\def\E{{\rm E}}
\def\Pvalue{{$p$-value}}
\def\Pvalues{{$p$-values}}
\def\pr{{\rm pr}}
\def\Pr{{\rm pr}}
\def\logit{{\rm logit}}
\def\T{{ \mathrm{\scriptscriptstyle T} }}
\def\diag{{\rm diag}}  
\def\half{{\textstyle{1\over 2}}}
\def\ith{{$i{\rm th}$ }}
\def\var{{\rm var}}
\def\cl{{{\rm c}\ell}}
\def\checkx{{\checked\hspace{-6pt}\setminus}}
\def\grad{{\nabla}}
\def\bs#1{{\boldsymbol{#1}}}

\newcommand{\R}{\mathbb{R}}
\newcommand{\N}{\mathbb N}
\newcommand{\Z}{\mathbb Z}
\newcommand{\Q}{\mathbb Q}
\newcommand{\1}{\mathds 1}
\newcommand{\Var}{\text{Var}}
\newcommand{\Cov}{\text{Cov}}
\newcommand{\sd}{\text{sd}}
\newcommand{\se}{\text{se}}
\renewcommand{\L}{\mathcal {L}}
\renewcommand{\lambda}{\uplambda}
\newcommand{\iidsim}{\stackrel{iid}{\sim}}
\newcommand{\otherwise}{\text{otherwise}}
\newcommand{\indep}{\perp \!\!\! \perp}
\renewcommand{\epsilon}{\varepsilon}
\newcommand{\st}{\text{ s.t. }}
\newcommand{\ds}{\displaystyle}
\newcommand{\ind}{\stackrel{d}{\to}}
\newcommand{\inp}{\stackrel{p}{\to}}
\newcommand{\inas}{\stackrel{a.s.}{\to}}
\DeclareMathOperator*{\argmax}{arg\,max}
\DeclareMathOperator*{\argmin}{arg\,min}

\sloppy %permits line-breaking of urls

\def\blue{\color{blue}}
\def\red{\color{red}}
\def\sgn{{\rm sign}}
\def\drightarrow{{\stackrel{d}{\rightarrow}}}
\begin{document}

\noindent{\bf Questions Week 7} {\hfill STA 2212S 2026}\\

%{\bf Due January 14} \hfill{\it MS, Exercise 5.2}

\bigskip

\begin{enumerate}

\item In Davison's analysis of publication bias,\footnote{Davison, A.C. {\it Statistical Models}, Chapter 5.5} he considers a model in which a single study with $n$ individuals leads to estimate $\hat\mu \dotsim N(\mu, \sigma^2/n)$; that this study is published, if some measure of randomness in publication, denoted $Z$, is positive, 
	and that $\hat\mu$ and $Z$ are related according to the model  $$\hat\mu = \mu + \sigma n^{-1/2} U_1, \quad Z = \gamma_0 + \gamma_1n^{1/2} + U_2, \quad \text{cor}(U_1,U_2)=\rho \ge 0.$$ 
	\begin{enumerate}
	\item Show that $\pr(R=1) = \pr(Z > 0) = \Phi(\gamma_0+\gamma_1n^{1/2})$ and 
	$$\displaystyle{\pr(R=1\mid\hat\mu) = \pr(Z > 0 \mid \hat\mu) = \Phi\left\{\frac{\gamma_0 + \gamma_1n^{1/2} + \rho n^{1/2}(\hat\mu-\mu)/\sigma)}{(1-\rho^2)^{1/2}}\right\}}.$$  
	\item  Show that the estimate of $\mu$ is biased if $\rho>0$: \begin{align*}E(\hat\mu \mid R=1) &= \mu+\rho\sigma n^{-1/2}\zeta(\gamma_0+\gamma_1 n^{1/2}), %\mu+\rho\sigma\gamma_1\zeta'(\gamma_0)+\rho\sigma\zeta(\gamma_0)n^{-1/2}
	\end{align*} where $\zeta(x) = \phi(x)/\Phi(x)$, and $\phi(\cdot), \Phi(\cdot)$ are the standard normal density and distribution function, respectively
	\item For a collection of $k$ such studies, each producing estimate $\hat\mu_j \dotsim N(\mu, \sigma^2/n_j)$, derive the log-likelihood function given in Davison's eq.(5.34): $$ \ell(\mu,\rho, \sigma^2) = \sum_{j=1}^k \left\{-\frac{1}{2}\log\sigma^2 + \frac{n_j}{2\sigma^2}(\hat\mu-\mu)^2 - \log\Phi(a_j)+\log\Phi(b_j) \right\}, $$
	where $a_j = \gamma_0+\gamma_1n_j^{1/2}$ and $ b_j = \left\{a_j+\rho n_j^{1/2}(\hat\mu_j-\mu)/\sigma\right\}(1-\rho^2)^{-1/2}$.
	\item Verify that if $\rho=0$ that the maximum likelihood estimate of $\mu$ is $$\frac{\sum_{j=1}^k n_j\hat\mu_j}{\sum_{j=1}^k n_j}.$$
\end{enumerate}

\textcolor{red}{
(a) By definition, 
\begin{align*}
	P(R = 1) &= P(Z > 0)\\
	&= P(U_2 > -\gamma_0 - \gamma_1 n^{1/2})\\
	&= 1 - P(U_2 \leq -\gamma_0 - \gamma_1 n^{1/2})\\
	&= \Phi(\gamma_0 + \gamma_1n^{1/2})
\end{align*}
By definition, $\hat\mu \sim \mathcal N(\mu, \sigma^2/n)$ and $Z \sim \mathcal N(\gamma_0 + \gamma_1n^{1/2}, 1)$. This implies 
\begin{equation*}
	\begin{pmatrix}
		Z \\ \hat\mu	
	\end{pmatrix} \sim \mathcal N_2\left(\begin{pmatrix}
		\gamma_0 + \gamma_1 n^{1/2} \\
		\mu
	\end{pmatrix}, \begin{pmatrix}
		1 & \rho\sigma n^{-1/2}\\
		\rho \sigma n^{-1/2} & \sigma^2 / n
	\end{pmatrix}\right)
\end{equation*}
so 
\begin{equation*}
	Z \mid \hat \mu \sim \mathcal N\left(\gamma_0 + \gamma_1 n^{1/2} +  \rho n^{1/2}\frac{\hat\mu - \mu}{\sigma}, 1 - \rho^2\right)
\end{equation*}
and hence 
\begin{align*}
	P(Z > 0 \mid \hat\mu) &= P\left(X > \frac{-\gamma_0 - \gamma_1n^{1/2}  - \rho n^{1/2}(\hat\mu - \mu) / \sigma}{(1 - \rho^2) ^{1/2} }\right) & \text{where $X \sim \mathcal N(0, 1)$}\\
	&= \Phi\left(\frac{\gamma_0 + \gamma_1 n^{1/2} + \rho n^{1/2} (\hat\mu - \mu)/\sigma}{(1 - \rho^2)^{1/2}}\right)
\end{align*}
as required.
}

\textcolor{red}{
(b) By definition, 
\begin{equation*}
	E[\hat\mu \mid Z > 0] = \mu + \sigma n^{-1/2} E[U_1 \mid U_2 > -\gamma_0 - \gamma_1n^{1/2}]
\end{equation*}
We know that 
\begin{equation*}
	\begin{pmatrix}
		U_1 \\ U_2
	\end{pmatrix} \sim \mathcal N_2\left(\begin{pmatrix}
		0 \\ 0
	\end{pmatrix}, \begin{pmatrix}
		1 & \rho \\ \rho & 1
	\end{pmatrix}\right)
\end{equation*}
By the law of iterated expectation, 
\begin{equation*}
	E[U_1 \mid U_2 > -\gamma_0 - \gamma_1 n^{1/2}] = E[E[U_1 \mid U_2] \mid U_2 > -\gamma_0 - \gamma_1 n^{1/2}]
\end{equation*}
Since $E[U_1 \mid U_2] = \rho U_2$, then 
\begin{align*}
	E[U_1 \mid U_2 >  -\gamma_0 - \gamma_1 n^{1/2}] &= \rho E[U_2 \mid U_2 >  -\gamma_0 - \gamma_1 n^{1/2}] \\
	&= \frac{E[U_2 \1_{U_2 >  -\gamma_0 - \gamma_1 n^{1/2}}]}{P(U_2 >  -\gamma_0 - \gamma_1 n^{1/2})}\\
	&= \frac{1}{\Phi(\gamma_0 + \gamma_1 n^{1/2})} \int_{ -\gamma_0 - \gamma_1 n^{1/2}}^{\infty} u\phi(u) du\\
	&= \frac{\phi(\gamma_0 + \gamma_1 n^{1/2})}{\Phi( \gamma_0 + \gamma_1 n^{1/2})}
\end{align*}
hence 
\begin{equation*}
	E(\hat \mu \mid R = 1) = \mu + \rho\sigma n^{-1/2} \zeta(\gamma_0 + \gamma_1 n^{1/2})
\end{equation*}
}

\textcolor{red}{
(c) Since we only observe the publisehd studies, the likelihood function is 
\begin{align*}
	\mathcal L(\mu, \rho, \sigma^2) &= \prod_{j=1}^k f(\hat \mu_j  \mid R_j = 1)\\
	&= \prod_{j=1}^k \frac{f(\hat\mu_j) P(R_j = 1 \mid \mu_j)}{P(R_j = 1)}\\\
	&= \prod_{j=1}^k \frac{\exp\left(-\frac{n_j}{2\sigma^2}(\hat\mu_j - \mu)^2\right)\Phi\left(\frac{\gamma_0 + \gamma_1 n^{1/2} + \rho n^{1/2} (\hat\mu - \mu)/\sigma}{(1 - \rho^2)^{1/2}}\right)}{\Phi(\gamma_0 + \gamma_1 n^{1/2})\sqrt{2\pi \sigma^2/n_j}}
\end{align*}
thus the log-likelihood is
\begin{equation*}
	\ell(\mu, \rho, \sigma^2) = \sum_{i=1}^n \left[-\frac 12 \log \sigma^2 - \frac{n_j}{2\sigma^2}(\hat\mu_j - \mu)^2 + \log \Phi(b_j) - \log \Phi(a_j)\right] + C
\end{equation*}
}

\textcolor{red}{
(d) When $\rho = 0$, 
\begin{equation*}
	\ell(\mu, 0, \sigma^2) = \sum_{j=1}^k \left[-\frac 12 \log \sigma^2 + \frac{n_j}{2\sigma^2}(\hat\mu_j - \mu)^2\right]
\end{equation*}
Differentiating wrt $\mu$ yields 
\begin{equation*}
	\frac{\partial \ell}{\partial \mu} = \sum_{j=1}^k \frac{n_j}{\sigma^2}(\mu - \hat\mu_j) = 0
\end{equation*}
Solving yields an MLE for $\mu$ of 
\begin{equation*}
	\hat\mu = \frac{\sum_{j=1}^k n_j \hat\mu_j}{\sum_{j=1}^k n_j}
\end{equation*}
as required. 
}

\item{\it Davison, Exercise 5.8, Applying the EM algorithm to censored data}
\begin{enumerate}
	\item Suppose $U$ is an underlying failure time, and that $Y=\min(U, c)$ and $D = I(U\le c)$ are observed, where $c$ is a fixed censoring time. The complete data log-likelihood function is $f(u;\theta)$. Show that $$f(u\mid y,d;\theta) = \begin{cases} \delta(u-y),&d=1, \\ \frac{f(u;\theta)}{1-F(c;\theta)}, & u>c, d=0, \end{cases}$$ where $\delta(u-y)$ is  the generalized derivative of the Heaviside function $H(u) = 1\{u\ge 0\}$, and satisfies $\int \delta(y-u)g(u)du = g(y)$ for any function $g(\cdot)$.\footnote{$\delta(\cdot)$ is called the Dirac delta function.}
	\item If $f(u;\theta) = \theta e^{-\theta u}$, show that $\E_{\theta'}(U\mid Y=y, D=d) = dy +(1-d)(c+1/\theta')$, and deduce that the estimate of $\theta$ at the $j$th iterate of the EM algorithm is $$\theta^{(j)} = \frac{n}{\sum_{j=1}^n\{d_jy_j + (1-d_j)(c+1/\theta^{ (j-1)})\}}.$$
	\item The basis of the EM algorithm is summarized by Eq.(5.36) in Davison: $$\log f(y,u;\theta) = \log f(y;\theta) + \log f(u\mid y;\theta);$$ here $U$ are the unobserved (latent) variables and the LHS is called the complete-data likelihood function.   Derive the equality 
    \begin{align*}
         -\frac{\partial^2\log f(y;\theta)}{\partial\theta\partial\theta^T} = \frac{\partial^2\ell(\theta)}{\partial\theta\partial\theta^T} &= \E_\theta\left\{-\frac{\partial^2\log f(y,U;\theta)}{\partial\theta\partial\theta^T}\mid Y=y \right\} \\ 
	    & ~~~~~~~~~~-\E_\theta\left\{-\frac{\partial^2\log f(U\mid y;\theta)}{\partial\theta\partial\theta^T}\mid Y=y\right\},
	\end{align*}
	i.e.~ $J(\theta) = I_c(\theta;y) - I_m(\theta;y)$, interpreted as meaning that the observed information equals the complete-data information minus the missing information.
	\item Show that the missing information for the model in (b) is $\sum_{j=1}^k(1-d_j)/\theta^2$. 
\end{enumerate}

\noindent \textcolor{red}{
(a) Suppose that $d = 1$. Then 
\begin{equation*}
	f(u \mid y, d = 1 ; \theta) = \delta(u-y)
\end{equation*}
since $u = y$ by definition. When $d = 0$ and $u > c$, 
\begin{equation*}
	f(u \mid y, d = 0 ; \theta) = \frac{P(Y = c, D = 0 \mid u)f(u)}{P(Y = c, D = 0)}
\end{equation*}
$D = 0$ implies that $U > c$, so $Y = c$, thus 
\begin{equation*}
	P(Y = c, d = 0) = P(U > c) = 1 - F(c ; \theta)
\end{equation*}
As well, when $u > c$, then $D = 0$ is guaranteed, so $P(Y = c, D = 0 \mid u) = 1$, 
hence 
\begin{equation*}
	f(u \mid y, d ; \theta) = \begin{cases}
		\delta(u - y) & d = 1\\
		\frac{f(u ; \theta)}{1 - F(c ; \theta)} & u > c , d = 0
	\end{cases}
\end{equation*}
}

\textcolor{red}{
(b) 
\begin{align*}
	E_{\theta'}[U \mid Y = y, D = d] &= d E_{\theta'}[U \mid Y = y, D = 1] + (1 - d) E_{\theta '} [U \mid Y = y, D = 0]\\
	&= d y + (1 - d) E_{\theta'} [U \mid U > c]
\end{align*}
To compute $E_{\theta'}[U \mid U > c]$, 
\begin{align*}
	E_{\theta'} [U \mid U > c] &= \frac{E_{\theta'} (U \1_{U > c})}{P(U > c)}\\
	&= \frac{1}{\exp(-\theta' c)} \int_c^\infty \theta' u \exp(-\theta' u)du\\
	&= c + \frac 1 {\theta'}
\end{align*}
thus 
\begin{equation*}
	E_{\theta'}[U \mid Y = y, D = d] = dy + (1 - d) \left(c + \frac{1}{\theta'}\right)
\end{equation*}
For the expectation step, we want to find $\theta^{(j)}$ defined as  
\begin{equation*}
	\theta^{(j)} = \argmax_{\theta} Q(\theta \mid \theta^{(j-1)})
\end{equation*}
where 
\begin{align*}
	Q(\theta \mid \theta^{(j-1)}) &= E_{\theta^{(j-1)}} (\ell(\theta) \mid y_1, d_1,\ldots, y_n, d_n)\\
	&= E_{\theta^{(j-1)}} \left(n \log \theta - \theta \sum_{j=1}^n U_j \mid y_1, d_1,\ldots y_n, d_n\right)\\
	&= n\log \theta - \theta \sum_{j=1}^n E(U_j \mid y_j, d_j)\\
	&= n\log \theta - \theta \sum_{j=1}^n \left[d_j y_j + (1 - d_j) \left(c + \frac{1}{\theta^{(j-1)}}\right)\right]
\end{align*}
Differentiating $Q$ wrt $\theta$ yields 
\begin{equation*}
	\frac{\partial Q}{\partial \theta} = \frac n \theta - \sum_{j=1}^n \left[d_j y_j + (1 - d_j) \left( c + \frac{1}{\theta^{(j-1)}}\right)\right]
\end{equation*}
thus 
\begin{equation*}
	\theta^{(j)} = \frac{n}{\sum_{j=1}^n \left[d_j y_j + (1 - d_j) \left( c + \frac{1}{\theta^{(j-1)}}\right)\right]}
\end{equation*}
}

\textcolor{red}{
(c) From the identity given, 
\begin{equation*}
	\frac{\partial^2 \log f(y ; \theta)}{\partial \theta \partial \theta^\top} = \frac{\partial^2 \log f(y, U ; \theta)}{\partial \theta\partial^\top} - \frac{\partial^2 \log f(U \mid y ; \theta)}{\partial\theta\partial\theta^\top}
\end{equation*}
Then conditioning on $Y =y$, the conditional expectation of $\partial^2\log f(y ; \theta) / \partial\theta\partial\theta^\top$ is constant, so 
\begin{equation*}
	-\frac{\partial^2 \log f(y ; \theta)}{\partial \theta \partial \theta^\top} = E_\theta \left[- \frac{\partial^2 \log f(y, U ; \theta)}{\partial \theta \partial \theta^\top} \mid Y = y \right] - E_\theta\left[-\frac{\partial^2 \log f(U \mid y; \theta)}{\partial \theta \partial \theta^\top} \mid Y = y\right]
\end{equation*}
}

\textcolor{red}{
(d) From (c), the missing information is defined as 
\begin{equation*}
	I_m(\theta ; y) = E_\theta \left[-\frac{\partial^2 \log f(U \mid y ; \theta)}{\partial \theta^2}\mid Y = y\right]
\end{equation*}
where 
\begin{equation*}
	f(u \mid y ; \theta) = \prod_{j=1}^k f(u_j \mid y_j, d_j ; \theta)
\end{equation*}
From part (a), if $d_j = 1$, 
\begin{equation*}
	\log f(u_j \mid y_j, 1 ; \theta) = \log \delta(u_j - y_j)
\end{equation*}
hence the derivatives with resepect to $\theta$ are $0$ in this case. For $d_j = 0$, 
\begin{equation*}
	\log f(u_j \mid y_j , 0 ; \theta) = \log f(u_j ; \theta) - \log(1 - F(c ; \theta)) = \log \theta - \theta u_j + \theta c
\end{equation*}
Taking derivatives twice gives 
\begin{equation*}
	\frac{\partial^2 \log f(u _j\mid y_j, 0 ; \theta)}{\partial\theta^2} = -\frac{1}{\theta^2}
\end{equation*}
So, the derivative is 
\begin{equation*}
	\frac{\partial^2 \log f(u_j \mid y_j, d_j; \theta)}{\partial\theta^2} = \begin{cases}
		0 & d _j = 1\\
		-\frac{1}{\theta^2} & d_j = 0
	\end{cases} = \frac{1 - d_j}{\theta^2}
\end{equation*}
Then the missing information is 
\begin{equation*}
	I_m(\theta ; y) = \sum_{j=1}^k E_\theta \left[-\frac{\partial \log f(U_j \mid y_j, d_j; \theta)}{\partial \theta^2}\mid Y_j = y_j\right] = \sum_{j=1}^k \frac{1-d_j}{\theta^2}
\end{equation*}
}
\end{enumerate}
\end{document}

